%=================================================
% Economics Coursework Template
%=================================================
\documentclass[12pt]{article}

%------------ Packages ------------
\usepackage[margin=1in]{geometry}   % Page layout
\usepackage{lmodern}                % Latin Modern (LaTeX default font)
\usepackage{amsmath, amssymb}       % Math symbols
\usepackage{booktabs}               % Professional tables
\usepackage{graphicx}               % Figures
\usepackage{natbib}                 % References
\usepackage[hidelinks]{hyperref}    % Hyperlinks
\usepackage{fancyhdr}               % Headers/footers
\usepackage{setspace}               % Line spacing
\usepackage{titlesec}               % Section formatting
\usepackage{geometry}
\usepackage{comment} 
\usepackage{xcolor} 
\usepackage{amsthm} 
\usepackage{float} 
\usepackage{dcolumn}
\usepackage{titlesec}
\usepackage{enumitem}
\usepackage{hyperref}
\usepackage{tikz}
\usepackage{multicol}
\usetikzlibrary{arrows.meta,positioning,fit,calc}
\usepackage{subcaption}
\tikzset{
  var/.style={rectangle, draw, rounded corners, minimum width=18mm, minimum height=7mm, align=center},
  latent/.style={ellipse, draw, dashed, minimum width=18mm, minimum height=7mm, align=center},
  controlarrow/.style={->, dashed, draw=red!60},
  lightbluearrow/.style={->, dashed, draw=blue!40} % NEW style
}

%------------ Page Style ------------
\rhead{\thepage}
\lhead{ECO1465 Weekly Report 1}

%------------ Section Formatting ------------
\titleformat{\section}{\large\bfseries}{\thesection.}{0.5em}{}
\titleformat{\subsection}{\normalsize\bfseries\itshape}{\thesubsection}{0.5em}{}
\titlespacing*{\subsection}{0pt}{1.5ex plus 1ex minus .2ex}{1em} % spacing before/after subsections


%------------ Begin Document ------------
\begin{document}
% Title Page only
\begin{titlepage}
    \thispagestyle{empty} % no headers/footers
    \centering
    \vspace*{\fill}

    {\Large \textbf{BLP} \par}

    \vspace{3em}
    {\large Nardeen Abdulkareem \par}
    {\normalsize The University of Toronto \par}
    {\normalsize \today \par}
    {\texttt{nardeen.abdulkareem@mail.utoronto.ca} \par}

    \vspace*{\fill}
\end{titlepage}
\clearpage

\section{Mis-specified Models}
\begin{table}[!htbp]\centering
\caption{Logit Demand Estimation Results}
\label{tab:demand_results}
\begin{threeparttable}
\footnotesize
\begin{tabular}{lcc}
\toprule
 & (1) OLS & (2) IV (2SLS) \\
\midrule
Price ($\alpha$) 
    & -0.804*** 
    & -1.791*** \\
    & (0.054) 
    & (0.069) \\
\\
$x$ 
    & 0.826*** 
    & 0.946*** \\
    & (0.029) 
    & (0.033) \\
\\
Satellite
    & 0.950*** 
    & 3.558*** \\
    & (0.141) 
    & (0.187) \\
\\
Wired
    & 0.921*** 
    & 3.512*** \\
    & (0.140) 
    & (0.184) \\
\midrule
Observations & 2400 & 2400 \\
$R^2$ & 0.286 & 0.144 \\
Estimator & OLS & IV-2SLS \\
\bottomrule
\end{tabular}
\begin{tablenotes}
\item Instrumented for price using cost shifter $w$. Robust standard errors. *** $p<0.01$, ** $p<0.05$, * $p<0.1$.
\end{tablenotes}

\end{threeparttable}
\end{table}

The OLS and IV estimates for the parameter of interest vary wildly. This is particularly clear with the price coefficient. Under OLS the price coefficient $\alpha$ is far larger than the true value of the price parameter and the estimated value using 2SLS. This is due to a endogenous price which present in the OLS model. Instrumenting for price using a cost shifter such as $\omega$ alleviates the endogenous price present when estimating parameters using OLS. The estimates for Satellite and Wired also increase, since price endogeneity influences how variation in prices and product characteristics is attributed to consumer utility. Once we instrument for prices, the model is able to capture more variation in demand which is attributed to product characteristics rather than price. This correctly implies that the 2SLS model reflects that consumers have much high sensitivity towards prices and product characteristics. 

\section{Correctly Specified Models}
\vspace{-0.1in}
\begin{table}[!htbp]\centering
\caption{Firm-Level Own-Price Elasticities}
\begin{threeparttable}
\footnotesize
\begin{tabular}{lcc}
\toprule
Firm & Estimated Elasticity & True Elasticity \\
\midrule
0 & -3.834 & -4.139 \\
1 & -3.881 & -4.184 \\
2 & -4.074 & -4.142 \\
3 & -4.088 & -4.160 \\
\bottomrule
\end{tabular}
\begin{tablenotes}
\footnotesize
\item Estimated elasticities of joint S \& D model. True elasticities of parameters from simulated data.
\end{tablenotes}

\end{threeparttable}
\end{table}

My preferred specification is the jointly estimated demand and supply model since it incorporates both the demand side and supply equations simultaneously into the equilibrium estimation. This model improves efficiency in the estimation of the parameters and yields information related to the supply sided measurements. The model also account for an endogenous price using the supply sided shifters, rival firms product characteristics, and rival firm cost shifters for the estimation of demand. Supply sided instruments used in the estimation include, own and rival firms demand shifter or product attributes. The joint model yields estimated elasticities which are very close to the true elasticities used to generate the simulated data. Although the estimated elasticities are slightly smaller in magnitude, they closely match the true values and correctly capture the overall level of price sensitivity in the data.

\begin{table}[H]
\centering
\caption{BLP Estimation Results}
\begin{tabular}{lcccc}
\toprule
 & Demand & Optimal Demand & Joint S \& D & Optimal Joint S \& D \\
\midrule
\textit{Demand parameters} \\
Price ($\alpha$) 
    & -1.8371 & -1.8786 & -1.8794 & -1.8706 \\
    & (2.5111) & (0.0765) & (0.1148) & (0.0849) \\
x
    & 0.9716 & 0.9929 & 0.9928 & 0.9919 \\
    & (1.4623) & (0.0367) & (0.0537) & (0.0371) \\
Satellite
    & 3.5620 & 3.5811 & 3.5844 & 3.5747 \\
    & (0.3228) & (0.1982) & (0.2116) & (0.2179) \\
Wired
    & 3.6178 & 3.7130 & 3.7153 & 3.6923 \\
    & (5.6523) & (0.2014) & (0.2793) & (0.2216) \\
\addlinespace
\textit{Supply parameters} \\
Constant
    & --- & --- & 0.4816 & 0.4772 \\
    & --- & --- & (0.0180) & (0.0229) \\
w
    & --- & --- & 0.2492 & 0.2558 \\
    & --- & --- & (0.0073) & (0.0318) \\
\midrule
\multicolumn{5}{l}{\footnotesize Robust standard errors in parentheses.} \\
\bottomrule
\end{tabular}
\end{table}

\section{Merger}
In our single product firm framework, each firm maximizes profits by internalizing the prices of rival firms. Firms compete under Bertrand competition. As such, the FOC condition for a firms profit maximization problem indicates that equilibrium prices are determined by the markup a firm can set on its product, the elasticity of demand for the good, and the substitution pattern across products. Prices are strategic complements because when one firm raises its price, rival firms face less competitive pressure and the demand they face becomes less elastic. As a result, the rival firm will be able to raise its prices when a rival firm increases its price. In theory, the magnitude of a firms response to a rivals price increase is determined by the proximity of a firm to a rival with a close substitute product. An high degree of substitutability between two product will yield a larger increase in price of their own good in response to an increase to the price of a rival firms good. This is because in a price setting environment firms competed strongly with each other, so each firm had an incentive to keep prices relatively low to avoid losing customers to the other.

If firms 1 and 2 merge, then the price of these goods rise since goods 1 and 2 no longer need to compete with one another for a share of the market. In this case, some consumers may switch to products 3 and 4 (wired services) so firms 3 and 4 can also raise their prices slightly to capture larger profits. Since these goods are close substitutes, the increase in price (for goods 1 and 2) observed will be larger than the increase in price for goods 1 and 3 if firms 1 and 3 merged. This is captured bellow in our model. 

\begin{table}[h!]
\centering
\caption{Comparison of Average Price Changes Across Merger Scenarios}
\begin{tabular}{lcccc}
\hline
Product & Baseline P & $\Delta P$ (Merger 1--2) & $\Delta P$ (Merger 1--3) & $\Delta P$ (Merger 1--2) w/ Efficiency\\
\hline
1 & 2.7487 & 0.2613 & 0.1235 & 0.0612 \\
2 & 2.7300 & 0.2744 & 0.0126 & 0.0751 \\
3 & 2.7176 & 0.0093 & 0.1312 & -0.0001 \\
4 & 2.7312 & 0.0098 & 0.0083 & 0.0000 \\
\hline
\end{tabular}
\end{table}

When firms 1 and 2 merge, prices of goods 1 and 2 increase substantially by around 9 to 10\%, while prices of goods 3 and 4 change very little. This is attributed to high substitutability between products 1 and 2. More interestingly, when the merger occurs between firms 1 and 3, where these two firms provide products which are not close substitutes to one another. The change in prices are more modest, at around 4 to 5\%, when compared to the change in prices due to the merger of firms 1 and 2. This indicates that goods 1 and 3 are weaker substitutes than goods 1 and 2.

So far we have no explored a reduction in the marginal cost associate with producing a good due to a merger. The efficiencies that occur as a result of a merger can be attributed to factors such as improved technology, economies of scale, or reduced redundancies. This increase in efficiency is unlikely to materialize as welfare enhancing on the consumer side. Firms benefit from lower production costs and potentially higher profits. Since total welfare equals consumer surplus plus producer surplus, the reduction in production costs represents a real efficiency gain that on the producer side. A large marginal cost reductions can lead to higher total welfare through increased producer surplus which would make the merger welfare enhancing.

\begin{table}[h!]
\centering
\caption{Welfare Effects of the Merger Between Firms 1 and 2 with a 15\% Cost Reduction}
\begin{tabular}{l c}
\hline
 & \textbf{Change} \\
\hline
Consumer Surplus & -0.0132 \\
Producer Profits & 0.1118 \\
Total Welfare & 0.0986 \\
\hline
\end{tabular}
\end{table}

The model predicts that with increased efficiency due to a merger between firms 1 and 2, prices of goods 1 and 2 rise by an amount substantially smaller had there not been any increased efficiency. Goods 3 and 4 do not exhibit any changes in prices in this case. The increase in the prices of goods 1 and 2 influences the consumer surplus negatively since a price increase hurts consumers when there is not adequate substitute product which they can chose to purchase instead. Regardless, a reduction in marginal cost can cause a merger to be welfare enhancing because it allows the firm to produce the good more efficiently. This increase in efficiency offsets the increase in mark ups caused by internalized substitution between the merged products. Since we observe both an increase in prices and a reduction in costs, the firm enjoys a larger producer surplus which outweighs the negative change in consumer surplus. Therefore, since the efficiency gains are large enough, the merger increases total welfare even if it increases market power.

The assumption made regarding $M_t$ is needed since measurement of welfare are dependant on the size of the market or the number of consumers in each market. Prices are determined in the first order conditions by estimating market shares which are normalized to the size of the market which causes markups to not depend on the number of consumers in the market. However, when we aggregate over consumers in order to attain welfare figures, the size of a market determines how salient the impact of a merger may be. For the sake of keeping a parsimonious analysis we fix $M_t = 1$

\section{Some Plots}

\begin{figure}[H]
\centering

% Top row
\begin{subfigure}{0.33\textwidth}
    \centering
    \includegraphics[width=\linewidth]{1.png}
\end{subfigure}
\hfill
\begin{subfigure}{0.45\textwidth}
    \centering
    \includegraphics[width=\linewidth]{2.png}
\end{subfigure}

\vspace{0.5cm}

% Bottom centered plot
\begin{subfigure}{0.33\textwidth}
    \centering
    \includegraphics[width=\linewidth]{3.png}
\end{subfigure}

\caption{Comparison of model results}
\label{fig:threeplots}
\end{figure}

\newpage
\section{Example of Strategic Relationships from Nielson, 2021}
Schools choose price $p_j$ and quality $q_j$ to maximize profits. The profit function of school $j$ is:
\begin{equation}
\pi_j(q,p,\xi) = \left(p_j + v_b^m - MC(q_j)\right) N_m s_j(q,p,\xi) - F_j
\end{equation}

Where $s_j(q,p,\xi)$ is the market share of school $j$, $v_b^m$ is the voucher per student in market $m$, $MC(q_j)$ is the marginal cost of quality, and $F_j$ is a fixed cost. We have the following properties for shares of students of type k, located at loc, who chooses school j. This aid in our subsequent analysis of strategic relationships.
\begin{equation}
s_{jk}^{loc}(q,op,\xi)
=
\frac{
\exp\!\left(
x_{jt}\beta
+
\beta_i q_{jt}
-
\alpha_i op_{kjt}
+
\lambda_i d_{loc,j}
+
\xi_{jt}
\right)
}{
\sum_{r}
\exp\!\left(
x_{rt}\beta
+
\beta_i q_{rt}
-
\alpha_i op_{krt}
+
\lambda_i d_{loc,r}
+
\xi_{rt}
\right)
}.
\end{equation}

\[
\frac{\partial s_{jk}}{\partial p_j} < 0, \quad
\frac{\partial s_{jk}}{\partial q_j} > 0, \quad
\frac{\partial s_{jk}}{\partial p_k} > 0, \quad
\frac{\partial s_{jk}}{\partial q_k} < 0.
\]

This occurs since we use an additive discrete choice model (logit) in the model. Where the household utility for household $i$ obtained from attending school $j$ in market $m$ and period $t$ is given by:


\begin{equation}
u_{ijt}
=
x_{jt}\beta
+
\beta_i q_{jt}
-
\alpha_i \, op_{k(i)jt}
+
\lambda_i d_{loc(i),j}
+
\xi_{jt}
+
\epsilon_{ijt},
\end{equation}

When we aggregate our shares across student types and locations we are able to attain the following relationship for the partials of market share:
\begin{equation}
s_j(q,op,\xi)
=
\sum_{k} \Pi_k^m
\sum_{loc \in L_m}
w_k^{loc}
\, s_{jk}^{loc}(q,op,\xi),
\end{equation}

\subsubsection*{Own price}

\begin{equation}
\frac{\partial s_j}{\partial p_j}
=
\sum_k \Pi_k^m
\sum_{loc}
w_k^{loc}
\frac{\partial s_{jk}^{loc}}{\partial p_j}
< 0.
\end{equation}

\subsubsection*{Own quality}

\begin{equation}
\frac{\partial s_j}{\partial q_j}
=
\sum_k \Pi_k^m
\sum_{loc}
w_k^{loc}
\frac{\partial s_{jk}^{loc}}{\partial q_j}
> 0.
\end{equation}

\subsubsection*{Rival price}

\begin{equation}
\frac{\partial s_j}{\partial p_k}
=
\sum_k \Pi_k^m
\sum_{loc}
w_k^{loc}
\frac{\partial s_{jk}^{loc}}{\partial p_k}
> 0.
\end{equation}

\subsubsection*{Rival quality}

\begin{equation}
\frac{\partial s_j}{\partial q_k}
=
\sum_k \Pi_k^m
\sum_{loc}
w_k^{loc}
\frac{\partial s_{jk}^{loc}}{\partial q_k}
< 0.
\end{equation}

In addition to the properties that we have laid down regarding market share above. We must assume the following;

\begin{enumerate}

\item Marginal cost of quality is increasing:
\[
MC'(q_j) > 0.
\]

\item The second-order condition for profit maximization holds:
\[
\frac{\partial^2 \pi_j}{\partial a_j^2} < 0.
\]

\item Direct demand effects dominate higher-order cross effects, so that the signs
of the first-order demand derivatives determine the sign of the profit cross-partials.
\end{enumerate}

\subsubsection*{Profit Maximization FOC}
Assuming $MC(q_j)=c_0+c_1 q_j$

Differentiate profits with respect to $p_j$:
\begin{align}
s_j(q,op,\xi)
+
\left(p_j^*+v_b^m-MC(q_j^*)\right)\frac{\partial s_j(q,op,\xi)}{\partial p_j}
&=0.
\end{align}

\begin{equation}
p_j^*
=
c_0+c_1 q_j-v_b^m
-
s_j(q,op,\xi)
\left[\frac{\partial s_j(q,op,\xi)}{\partial p_j}\right]^{-1}.
\end{equation}

Differentiate profits with respect to $q_j$:
\begin{align}
-\,MC'(q_j^*)\,s_j(q,op,\xi)
+
\left(p_j^*+v_b^m-MC(q_j^*)\right)\frac{\partial s_j(q,op,\xi)}{\partial q_j}
&=0.
\end{align}


\begin{equation}
q_j^*
=
\frac{p_j^*+v_b^m-c_0}{c_1}
-
s_j(q,op,\xi)
\left[\frac{\partial s_j(q,op,\xi)}{\partial q_j}\right]^{-1}.
\end{equation}

\subsubsection*{1. Own price and own quality: $(p_j,q_j)$}

From
\[
\frac{\partial \pi_j}{\partial p_j}
=
N_m\left[
s_j + A_j\frac{\partial s_j}{\partial p_j}
\right],
\]
differentiate with respect to $q_j$:

\begin{align}
\frac{\partial^2 \pi_j}{\partial q_j \partial p_j}
&=
N_m\left[
\frac{\partial s_j}{\partial q_j}
+
\frac{\partial A_j}{\partial q_j}\frac{\partial s_j}{\partial p_j}
+
A_j\frac{\partial^2 s_j}{\partial q_j \partial p_j}
\right] \\
&=
N_m\left[
\frac{\partial s_j}{\partial q_j}
-
MC'(q_j)\frac{\partial s_j}{\partial p_j}
+
A_j\frac{\partial^2 s_j}{\partial q_j \partial p_j}
\right].
\end{align}

Since $\frac{\partial s_j}{\partial q_j}>0$ and $\frac{\partial s_j}{\partial p_j}<0$, the first two terms are positive when $MC'(q_j)>0$. Hence
\[
\frac{\partial^2 \pi_j}{\partial q_j \partial p_j}>0.
\]

Therefore, own price and own quality are \textbf{strategic complements}.

\subsubsection*{2. Own price and rival price: $(p_j,p_k)$}

Differentiate $\partial \pi_j / \partial p_j$ with respect to $p_k$:

\begin{align}
\frac{\partial^2 \pi_j}{\partial p_k \partial p_j}
&=
N_m\left[
\frac{\partial s_j}{\partial p_k}
+
A_j\frac{\partial^2 s_j}{\partial p_k \partial p_j}
\right].
\end{align}

Since $\frac{\partial s_j}{\partial p_k}>0$, the cross-derivative is positive under the standard assumption that the direct demand-shifting effect dominates:
\[
\frac{\partial^2 \pi_j}{\partial p_k \partial p_j}>0.
\]

Therefore, prices are \textbf{strategic complements}.

\subsubsection*{3. Own quality and rival quality: $(q_j,q_k)$}

From
\[
\frac{\partial \pi_j}{\partial q_j}
=
N_m\left[
A_j\frac{\partial s_j}{\partial q_j}
-
MC'(q_j)s_j
\right],
\]
differentiate with respect to $q_k$:

\begin{align}
\frac{\partial^2 \pi_j}{\partial q_k \partial q_j}
&=
N_m\left[
A_j\frac{\partial^2 s_j}{\partial q_k \partial q_j}
-
MC'(q_j)\frac{\partial s_j}{\partial q_k}
\right].
\end{align}

Since $\frac{\partial s_j}{\partial q_k}<0$, the second term is positive. Hence
\[
\frac{\partial^2 \pi_j}{\partial q_k \partial q_j}>0.
\]

Therefore, qualities are \textbf{strategic complements}.

\subsubsection*{4. Own price and rival quality: $(p_j,q_k)$}

Differentiate $\partial \pi_j / \partial p_j$ with respect to $q_k$:

\begin{align}
\frac{\partial^2 \pi_j}{\partial q_k \partial p_j}
&=
N_m\left[
\frac{\partial s_j}{\partial q_k}
+
A_j\frac{\partial^2 s_j}{\partial q_k \partial p_j}
\right].
\end{align}

Since $\frac{\partial s_j}{\partial q_k}<0$, the cross-derivative is negative under the standard assumption that the direct demand-loss effect dominates:
\[
\frac{\partial^2 \pi_j}{\partial q_k \partial p_j}<0.
\]

Therefore, own price and rival quality are \textbf{strategic substitutes}.

\subsubsection*{5. Own quality and rival price: $(q_j,p_k)$}

Differentiate $\partial \pi_j / \partial q_j$ with respect to $p_k$:

\begin{align}
\frac{\partial^2 \pi_j}{\partial p_k \partial q_j}
&=
N_m\left[
A_j\frac{\partial^2 s_j}{\partial p_k \partial q_j}
-
MC'(q_j)\frac{\partial s_j}{\partial p_k}
\right].
\end{align}

Since $\frac{\partial s_j}{\partial p_k}>0$, the second term is negative. Hence
\[
\frac{\partial^2 \pi_j}{\partial p_k \partial q_j}<0.
\]

Therefore, own quality and rival price are \textbf{strategic substitutes}.

\subsubsection*{Conclusion}

\[
\frac{\partial^2 \pi_j}{\partial q_j \partial p_j}>0
\qquad\Rightarrow\qquad
(p_j,q_j)\text{ are strategic complements;}
\]

\[
\frac{\partial^2 \pi_j}{\partial p_k \partial p_j}>0
\qquad\Rightarrow\qquad
(p_j,p_k)\text{ are strategic complements;}
\]

\[
\frac{\partial^2 \pi_j}{\partial q_k \partial q_j}>0
\qquad\Rightarrow\qquad
(q_j,q_k)\text{ are strategic complements;}
\]

\[
\frac{\partial^2 \pi_j}{\partial q_k \partial p_j}<0
\qquad\Rightarrow\qquad
(p_j,q_k)\text{ are strategic substitutes;}
\]

\[
\frac{\partial^2 \pi_j}{\partial p_k \partial q_j}<0
\qquad\Rightarrow\qquad
(q_j,p_k)\text{ are strategic substitutes.}
\]

\paragraph{(1) Own price and own quality: $(p_j,q_j)$}

Price and quality are strategic complements.

Higher quality increases the willingness to pay of families. As a result,
an increase in quality raises the profitability of increasing price. Conversely,
if the school charges a higher price, it has stronger incentives to provide higher
quality to justify the price and maintain enrollment.

\paragraph{(2) Prices: $(p_j,p_k)$}

Prices are strategic complements.

When a rival school increases its price, some students switch to school $j$,
which reduces the elasticity of demand faced by school $j$. With less elastic
demand, school $j$ can profitably increase its own price.

\paragraph{(3) Qualities: $(q_j,q_k)$}

Qualities are strategic complements.

If rival schools improve their quality, school $j$ loses demand unless it also
improves quality. Therefore, an increase in rival quality increases the incentive
for school $j$ to raise its own quality.

\paragraph{(4) Own price and rival quality: $(p_j,q_k)$}

Own price and rival quality are strategic substitutes.

If a rival school improves its quality, it becomes more attractive to students,
which reduces demand for school $j$. In response, school $j$ lowers its price to
remain competitive.

\paragraph{(5) Own quality and rival price: $(q_j,p_k)$}

Own quality and rival price are strategic substitutes.

If rival schools increase their prices, some students shift toward school $j$
even without improvements in quality. As a result, school $j$ faces weaker
incentives to invest in quality when rival prices increase.








\end{document}


